% =====================================================================
%  INTRODUCCIÓN
% =====================================================================
\chapter*{Introducción}
\addcontentsline{toc}{chapter}{Introducción}

El acoso sexual en las universidades es un problema conocido, pero poco denunciado. Las relaciones de jerarquía, la costumbre de restar importancia a ciertas conductas y el temor a las consecuencias hacen que muchas personas prefieran callar \citep{bondestam2020, rios2026}. La UNSAAC no es ajena a esta realidad: en una encuesta aplicada a sus estudiantes, el 77,5\,\% reportó haber vivido al menos una manifestación de hostigamiento sexual, y menos del 1\,\% de quienes lo vivieron presentó una denuncia \citep{gutierrez2022}. La Ley N.\textsuperscript{o}~27942 y su reglamento obligan a las universidades a prevenir y atender estos casos \citep{ley27942, ds0142019}, pero difícilmente se puede atender lo que no llega a conocerse.

Una parte de estas experiencias sí se cuenta, aunque no por la vía formal: aparece en testimonios, en conversaciones con tutores o en mensajes a las oficinas de orientación. Leer esos relatos y reconocer cuándo describen acoso no es sencillo, porque el significado de una situación depende del contexto y de la relación entre las personas, y muchas veces no hay una sola palabra que lo delate \citep{salawu2020, dutta2025}.

La forma más intuitiva de automatizar esa lectura sería buscar palabras clave, pero ese enfoque falla justo donde más importa: no reconoce el acoso que se cuenta sin términos explícitos y marca como acoso cualquier texto que los mencione. El aprendizaje automático ofrece otra salida. Un modelo entrenado con ejemplos etiquetados aprende qué combinaciones de palabras y qué contextos acompañan a un relato de acoso, y puede aplicar ese aprendizaje a relatos nuevos \citep{salawu2020}. Los clasificadores clásicos, como la regresión logística o las máquinas de vectores de soporte, funcionan bien con pocos datos y son fáciles de interpretar; los modelos basados en \textit{transformers}, como BETO, representan el significado de cada palabra según la oración en que aparece \citep{devlin2019, canete2020}, aunque necesitan más cómputo. Por eso esta tesis no elige uno de antemano: los compara en igualdad de condiciones. Esta idea ya se ha probado en otros contextos. \citet{gorissen2023} identificaron revelaciones de victimización sexual en mensajes en neerlandés; \citet{wijaya2022} clasificaron relatos de acoso académico según su tipo, y \citet{mohammadi2024} detectaron sexismo en redes sociales con modelos \textit{transformers}. En español también hay avances: \citet{escobar2022} clasificaron automáticamente los testimonios de violencia sexual publicados con la etiqueta \#Cuéntalo, \citet{imbaquinga2025} detectaron mensajes de cibergrooming con un modelo integrado en un prototipo web, \citet{cumba2018} entrenó una red neuronal para detectar ciberacoso en tuits y, en el Perú, \citet{aguirre2022} aplicaron BERT para detectar ciberacoso en Twitter. Sin embargo, todos estos trabajos usan publicaciones de redes sociales o mensajes, no relatos de estudiantes universitarios, y algunos se evaluaron de una manera que puede inflar los resultados \citep{kapoor2023}. En los antecedentes revisados no se encontró un modelo entrenado con relatos de estudiantes universitarios peruanos para detectar acoso sexual.

Esta tesis propone desarrollar ese modelo. La idea es simple de enunciar: dado un relato escrito por un estudiante, el modelo debe indicar si describe o no una situación de acoso sexual. Para lograrlo se construirá un corpus de 400 a 600 relatos aportados voluntariamente por estudiantes de la UNSAAC, anonimizados y etiquetados con una guía revisada por especialistas. Con ese corpus se entrenarán y compararán varios modelos, y el elegido se probará una sola vez con relatos de participantes que no intervinieron en el entrenamiento. Su desempeño se medirá con el macro F1, la precisión y la exhaustividad, y se comparará con una línea base léxica que solo busca palabras clave. Finalmente, el modelo se integrará en un prototipo que sirva de apoyo a las personas encargadas de revisar estos casos.

En ese marco, la investigación responde al siguiente problema general: ¿en qué medida el \proc{desarrollo de} \tec{un modelo de aprendizaje automático} permite \ctx{la detección de acoso sexual en relatos de estudiantes de la UNSAAC}, en comparación con una línea base de referencia? El objetivo general es \proc{desarrollar} \tec{un modelo de aprendizaje automático} \ctx{para la detección de acoso sexual en relatos de estudiantes de la UNSAAC}, cuyo desempeño se evaluará mediante el macro F1 y la exhaustividad en relatos de prueba, frente a una línea base léxica. Para alcanzarlo se plantean cuatro objetivos específicos: \proc{construir} un corpus anonimizado y anotado de \ctx{relatos de estudiantes de la UNSAAC}; \proc{entrenar y seleccionar} el \tec{modelo} con mejor desempeño, con y sin tratamiento del desbalance de clases; \proc{validar} el \tec{modelo seleccionado} frente a una línea base y analizar sus errores; e \proc{implementar} un prototipo que aplique el \tec{modelo validado} a nuevos \ctx{relatos}. La hipótesis general sostiene que el \tec{modelo de aprendizaje automático} \proc{desarrollado} \ctx{detecta acoso sexual en relatos de estudiantes de la UNSAAC} con un desempeño significativamente superior al de una línea base de referencia.

El valor de este trabajo es de tres tipos. En lo científico, aporta evidencia sobre qué tan bien aprende un modelo a reconocer el acoso en relatos en español y en qué casos se equivoca. En lo metodológico, deja un procedimiento para recolectar, anonimizar, etiquetar y evaluar textos sensibles sin que los datos de prueba se filtren en el entrenamiento. En lo tecnológico, entrega un modelo validado y un prototipo que la universidad podría usar como apoyo en sus oficinas de orientación.

El documento se organiza en tres capítulos. El primero presenta los aspectos generales: el problema, los objetivos, la justificación, la delimitación, las hipótesis, las variables y el método, es decir, qué se investiga y cómo se comprobará. El segundo reúne el marco teórico: los antecedentes que muestran el vacío, las bases matemáticas de los modelos y las métricas, y las tecnologías que se usarán. El tercero describe cómo se desarrollará la propuesta en cuatro etapas, una por cada objetivo específico, desde la recolección de los relatos hasta el prototipo. Al final se incluyen las conclusiones, las recomendaciones, la bibliografía y los anexos.
